\begin{abstract}
Chicken is a simple model of escalation: if both players push, both can lose badly; if one pushes and the other yields, the pusher often gains at the other's expense. Chickenizer is a Python system for a turn-based variant with richer state (health, resilience, and preference weights). It combines three views of the same game: strategy rules encoded as logic and checked with SymPy~\cite{sympy}; a simulation engine that plays strategies round by round; and a $2\times 2$ normal form built from one-shot counterfactuals, with pure and mixed Nash equilibria from Nashpy where applicable~\cite{nashpy,nash1950}. A separate repeated-match layer summarizes long simulations (action counts, resilience over time, and simple conditional summaries by last round's joint move); it is simulation-based, not a full repeated-game Nash solution. Optional Streamlit panels~\cite{streamlit} and small Q-learning experiments add demos and benchmarks. Evidence is automated: \filepath{python -m pytest unit_tests integration_tests -q} reports 247 passing tests on the evaluated snapshot. The main payoff is a clear pipeline for comparing logic, equilibrium predictions, engine behavior, and learning runs under assumptions stated in code and tests.
\end{abstract}
